\chapter*{Introduction}
\addcontentsline{toc}{chapter}{Introduction}

Public-key cryptography is based on the idea that some mathematical operations
are easy to perform but difficult to reverse without secret information. This
principle makes it possible to communicate securely without first sharing a
private key through a secure channel. Since its introduction in the 1970s,
public-key cryptography has become one of the central parts of modern
information security. Many public-key systems are based on computational
problems that are believed to be hard, such as integer factorization, discrete
logarithms, or variants of lattice problems.

One of the early ideas for constructing a public-key cryptosystem was to use the
subset-sum problem, also known as the knapsack problem. In this problem, one is
given a sequence of positive integers and a target sum, and the task is to decide
whether some subset of the sequence adds up to the target. The problem is
important in complexity theory because its decision version is NP-complete. This
made it an attractive candidate for cryptography: a sender could encrypt a
message by forming a subset sum, while an attacker would appear to face a hard
combinatorial problem.

The first well-known cryptosystem based on this idea was proposed by Merkle and
Hellman~\cite{MerkleHellman1978}. Its construction uses a private
superincreasing sequence, for which the subset-sum problem can be solved
efficiently by a greedy algorithm. This easy private instance is transformed by
modular multiplication, and optionally by a permutation, into a public sequence
that is intended to look like a hard general subset-sum instance. The legitimate
receiver can reverse the transformation and solve the easy private problem,
while an attacker only sees the public knapsack.

This idea is simple and elegant, but its history also shows an important problem
in cryptographic design. The hardness of the general subset-sum problem does not
automatically imply the security of a cryptosystem built from it. A
cryptosystem does not use arbitrary subset-sum instances; it generates instances
with hidden structure so that the legitimate user can decrypt efficiently. If
this structure remains visible in the public key, then the system may be much
weaker than the general problem suggests.

This is exactly what happened to classical knapsack cryptosystems. The basic
Merkle--Hellman scheme was broken by Shamir, who showed that an equivalent
trapdoor can be recovered in polynomial time~\cite{Shamir1982}. Later work also
showed that several modified or iterated knapsack constructions remained
vulnerable, often through structural or lattice-based attacks. In particular,
low-density subset-sum instances can often be attacked by lattice reduction
methods, which makes the density of a knapsack instance an important parameter
when discussing practical security. These results changed the role of knapsack
cryptosystems: they are no longer considered suitable for practical public-key
encryption, but they remain useful for studying the relation between
computational hardness, trapdoor design, and cryptanalysis.

The aim of this thesis is to study the classical Merkle--Hellman knapsack
cryptosystem, explain its mathematical background, compare selected variations,
and implement several related constructions for experimentation. The thesis does
not try to propose a secure modern knapsack cryptosystem. Instead, it treats
knapsack cryptography as a historical and educational case study. The main focus
is on understanding how the original construction works, why its security
assumption was too optimistic, and how different modifications tried to repair
its weaknesses.

The first chapter introduces the theoretical background needed for the rest of
the thesis. It explains basic cryptographic terminology, the distinction between
symmetric and asymmetric cryptography, the role of computational hardness, and
the subset-sum problem. It then describes the Merkle--Hellman cryptosystem in
detail, including key generation, encryption, and decryption.

The second chapter discusses weaknesses and variations of knapsack
cryptosystems. It first explains why the basic Merkle--Hellman construction is
vulnerable, including Shamir's polynomial-time attack and low-density
lattice-based attacks. It then presents selected attempts to modify or improve
the original construction, such as permuted and iterated variants, and also
discusses the Chor--Rivest knapsack-type cryptosystem. The chapter emphasizes
the difference between the theoretical hardness of subset-sum and the concrete
security of structured public keys.

The third chapter describes the implementation used in this thesis. The program
implements the classical Merkle--Hellman scheme and selected variants in a
common framework. It supports key generation, encryption, decryption, benchmark
runs, and simple subset-sum solver experiments. The implementation is intended
mainly as an experimental tool for checking correctness, comparing running
times, and illustrating how parameter choices influence generated knapsack
instances.

The fourth chapter presents the experimental results. The experiments verify
that the implemented schemes correctly encrypt and decrypt tested inputs,
compare the running time of the implemented variants, and examine the influence
of key-generation parameters on knapsack density. Additional experiments compare
brute force and meet-in-the-middle subset-sum solvers on small instances. These
experiments are not security proofs, but they support the theoretical discussion
by showing the behaviour of the implemented schemes and the practical growth of
simple solving methods.

Overall, this thesis uses knapsack cryptosystems to illustrate a broader lesson
in cryptography: a hard mathematical problem is not enough by itself. The way
the problem is used, the structure of generated instances, and the available
cryptanalytic methods are equally important. The history of knapsack
cryptography is therefore a useful example of why cryptographic constructions
must be analyzed as concrete systems, not only as applications of hard problems.
